\FloatBarrier
\subsection{Limitations}
\label{subsec:limitations}
\rtag{E7, E4, R1.5, R1.6, R1.7(3)}\begin{revblock}The energy, channel and twin environments are simulated rather than measured, the evaluation uses one small tabular data set, and the sensitivity of ATEF-DT is slightly below that of full participation. The reputation layer fails when adversaries form half of a selected set or collude, the score remains a heuristic for the long-run problem, and the framework provides data locality but no formal privacy guarantee.\end{revblock}

\section{Conclusion}
\label{sec:conclusion}
This article presented ATEF-DT, a framework that schedules battery-constrained clients of a federated digital-twin healthcare system with a single weighted-product score over predicted energy, participation and reputation, a twin-driven battery screen, per-client adaptive quantization and reputation-weighted aggregation. The score was formulated as a weighted-product scalarization with three proved properties, and a convergence bound was given for the update that is applied, with a bound on the participation bias in terms of the effective aggregation weights. On the WDBC data with simulated devices and channels, the framework matched the accuracy of full participation within the variation across seeds, used \V{sav-rand}\% less device energy than uniform sampling at the same budget, and reached a Jain participation index of \V{jain-atef}, against \V{jain-eafl} for energy-greedy selection. With 20\% label-flipping clients it retained \V{lf20-atef}\% accuracy. The comparison with alternative schedulers showed that the product rule is as efficient as tuned additive rules and, unlike them, insensitive to the scale of its criteria; the robustness study showed that the reputation layer protects through graded weights within a limited domain and fails outside it.

Future work will evaluate the framework on hardware and on longitudinal clinical data with demographic attributes, replace the within-round reference of the credibility engine by one that does not depend on the selected cohort, and combine the scheduler with client-level differential privacy under the constraint that selection is deterministic.

\section*{Declaration of generative AI and AI-assisted technologies}
\rtag{E6, R2.6}\rev{Figure~\ref{fig:architecture} is a schematic drawn by the authors in draw.io from boxes, arrows and typeset equations, and all result figures are produced by the released plotting script from the logged data; the manuscript contains no image produced by a generative model. During the preparation of this revision the authors used a large language model (Claude, Anthropic) to assist with the restructuring and language editing of the text, with the implementation of the simulation code and with the drawing of Figure~\ref{fig:architecture}. The authors reviewed and verified the resulting text, code and results, and take full responsibility for the content of the article.}
